\bookchapter{Bab 5: Buat Telefon Lebih Memahami Anda — Tetapan Asas}

\chapterintro{Luangkan beberapa minit untuk tetapan, peningkatan pengalaman sangat besar. Tetapan ini seperti “menala” telefon, menjadikannya lebih sesuai dengan tabiat penggunaan anda.}

\booksection{Satu, Bunyi dan GetaranSatu, Bunyi dan Getaran}

: Bergantung pada getaran untuk peringatan semasa senyap.

\begin{itemize}
  \item \textbf{Mod jangan ganggu}: Boleh menetapkan pengecualian, contohnya membenarkan panggilan dari senarai kenalan.
  \item \textbf{Pengasingan volum}: Volum media, volum nada dering, volum penggera boleh dilaraskan secara berasingan.
  \item \textbf{Dua, Paparan dan Kecerahan}automatik
  \item \textbf{: Biar telefon melaraskan secara automatik mengikut cahaya persekitaran, menjimatkan kerisauan.}Mod pelindung mata / paparan malam
\end{itemize}

\booksection{Dua, Paparan dan KecerahanDua, Paparan dan Kecerahan}

\begin{itemize}
  \item \textbf{Mod gelap}: Menjimatkan kuasa (sangat berkesan untuk skrin OLED), dan tidak menyilaukan mata pada waktu malam. Banyak App juga menyokong mod gelap.
  \item \textbf{Kadar segar semula skrin}: Sebahagian telefon Android boleh dilaraskan, semakin tinggi (seperti 120Hz), semakin lancar tatalan.
  \item \textbf{Kunci skrin automatik}: Tetapkan 15 saat atau 30 saat, selamat dan menjimatkan kuasa.
  \item \textbf{Saiz fon}: Rasa fon kecil? Besarkan sedikit.
  \item \textbf{Tiga, Kertas Dinding, Tema dan Skrin Utama}Tukar kertas dinding yang cantik (statik atau dinamik), mood pun berubah. Telefon jenama Android sering mempunyai kedai tema, boleh menukar keseluruhan ikon dan gaya.
  \item \textbf{Widget skrin utama}: Cuaca, jadual, kawalan muzik diletakkan di skrin utama, akses satu langkah. iOS dan Android kedua-duanya menyokong, tekan lama kawasan kosong skrin utama untuk menambah.
\end{itemize}

\booksection{Tiga, Kertas Dinding, Tema dan Skrin UtamaTiga, Kertas Dinding, Tema dan Skrin Utama}

Bahagian ini adalah fungsi “selepas ditetapkan, tidak boleh kembali”, setiap satu boleh meningkatkan kecekapan harian anda dengan ketara.

\textbf{1. Mod Fokus / Jangan Ganggu: Kembalikan Perhatian yang Dirampas oleh Pemberitahuan}Idea teras

\booksection{Empat, Tetapan Kecekapan: Jadikan Telefon sebagai “Otak Luar” AndaEmpat, Tetapan Kecekapan: Jadikan Telefon sebagai “Otak Luar” Anda}

Cadangan penggunaan

\booksubsection{Mod Fokus / Jangan Ganggu: Kembalikan Perhatian yang Dirampas oleh Pemberitahuan}

\begin{itemize}
  \item \textbf{· Mod kerja: Hanya benarkan kenalan "bintang" dan App perlu melalui;}· Mod tidur: Diaktifkan secara automatik, hanya benarkan penggera;
  \item \textbf{· Kecergasan / Membaca / Mesyuarat: Disesuaikan mengikut aktiviti semasa.}\newline : Gunakan automasi + pencetus lokasi / masa, biar mod bertukar secara automatik. Contohnya: Pintasan iOS boleh menetapkan "tiba di pejabat secara automatik masuk mod kerja".\newline 2. Pengurusan Pemberitahuan: Hentikan Titik Merah dari "Panggilan Berterusan"\newline Matikan pemberitahuan pemasaran
  \item \textbf{: 90\% pemberitahuan sebenarnya tidak anda perlukan.}Susun mengikut kepentingan
\end{itemize}

\booksubsection{Pengurusan Pemberitahuan: Hentikan Titik Merah dari “Panggilan Berterusan”}

\begin{itemize}
  \item \textbf{Pemberitahuan berkumpulan}: Lipat pemberitahuan yang sama jenis bersama, mengelakkan ditenggelami oleh puluhan mesej selepas satu buka kunci.
  \item \textbf{3. Gerak isyarat pintas: Jadikan Operasi Lebih Pantas}Tangkap layar
  \item \textbf{: Kuasa + Volum bawah (iPhone skrin penuh) / Kuasa + Volum bawah (Android umum).}Tangkap layar panjang / Tangkap layar tatal
\end{itemize}

\booksubsection{Gerak isyarat pintas: Jadikan Operasi Lebih Pantas}

\begin{itemize}
  \item \textbf{Imbas kod satu klik}: Tambah ikon "Imbas" di pusat kawalan / tetapan pintas.
  \item \textbf{Ketuk belakang (iPhone)}: Dalam [Tetapan] > [Fungsi Pembantu] > [Sentuh] > [Ketuk Belakang], ikatkan tindakan seperti "tangkap layar" atau "kunci skrin".
  \item \textbf{Pengaktifan suara}: Siri / Xiao Ai Tongxue / Xiao Yi / Xiao Bu, panggil satu kali boleh menyemak cuaca, tetapkan penggera.
  \item \textbf{4. Kaedah Input: Alat Kecekapan Paling Dipandang Rendah}Pilih kaedah input yang betul
  \item \textbf{:}· Sogou / Baidu: Perbendaharaan kata besar, kemas kini bahasa popular pantas;
\end{itemize}

\booksubsection{Kaedah Input: Alat Kecekapan Paling Dipandang Rendah}

\begin{itemize}
  \item \textbf{· Gboard (Google): Ringkas, penyegerakan rentas platform kuat, menyokong input gelongsor;}\newline Suis yang wajib\newline : Pembetulan pintar, perkataan cadangan, penyegerakan perbendaharaan kata awan, input angka dan simbol dengan geser ke atas.\newline Peringatan privasi\newline : Berhati-hati menggunakan fungsi "cadangan ayat penuh" atau "cadangan peribadi", mungkin memuat naik kandungan input.
  \item \textbf{5. Skrin Utama dan Multitugas: Lancarkan Aliran Kerja Anda}Pembahagian skrin App / Gambar dalam gambar
  \item \textbf{: Balas mesej semasa menonton video, buat nota semasa melayari internet.}Tetingkap kecil / Tetingkap terapung
\end{itemize}

\booksubsection{Skrin Utama dan Multitugas: Lancarkan Aliran Kerja Anda}

\begin{itemize}
  \item \textbf{App kerap di atas}: Bar Dock (bahagian bawah skrin) hanya letak 4 App yang paling kerap digunakan, kurangkan masa mencari.
  \item \textbf{Senarai Semakan "Hari Pertama Bekerja"}Selepas mendapat telefon baharu atau kemas kini versi besar, disyorkan meluangkan 10 minit untuk menyemak tetapan ini:
  \item \textbf{(1) Kata laluan kunci skrin + pengesahan biometrik (2) Pemberitahuan berkumpulan + matikan pemberitahuan pemasaran (3) Saiz fon / kadar segar semula / mod gelap}(4) Kaedah input lalai / pelayar / e-mel (5) Kenalan kecemasan / kad perubatan (6) Tetapan sandaran
\end{itemize}

\begin{tcbinfo}{Lima, Fungsi Pembantu dan Aksesibiliti: Kuasa Teknologi untuk Kebaikan}
Jangan takut dengan nama “fungsi pembantu”, fungsi-fungsi ini boleh digunakan oleh semua orang. Pada asalnya direka untuk orang kurang upaya, tetapi juga sangat berguna untuk orang biasa — tangan sementara cedera atau memeluk bayi, fon telinga rosak di dalam kereta api dan tidak dapat mendengar dengan jelas, memakai cermin mata baca untuk melihat fon kecil... semuanya boleh menyelamatkan keadaan. Pada akhirnya, reka bentuk aksesibiliti mencerminkan konsep\newline teknologi untuk kebaikan\newline : membolehkan semua orang menikmati kehidupan digital secara sama rata.

\end{tcbinfo}

\booksection{Lima, Fungsi Pembantu dan Aksesibiliti: Kuasa Teknologi untuk KebaikanLima, Fungsi Pembantu dan Aksesibiliti: Kuasa Teknologi untuk Kebaikan}

Pembaca skrin adalah alat teras bagi pengguna celik untuk menggunakan telefon. Ia akan\textbf{membacakan semua kandungan yang kelihatan pada skrin}, pengguna mengendalikan telefon melalui gerak isyarat dan suara. Berikut adalah pengenalan kaedah penggunaan iOS “VoiceOver” dan Android TalkBack.

\booksubsection{“Mata” Pengguna Celik — Panduan Lengkap Pembaca Skrin}

Laluan pengaktifan\textbf{: [Tetapan] > [Fungsi Pembantu] > [VoiceOver], buka suis. Boleh juga tiga klik kekunci sisi (atau kekunci skrin utama) untuk buka/tutup dengan pantas.}Gerak isyarat asas

\begin{tcbinfo}{}
\textbf{Luncur kiri-kanan satu jari: Tukar elemen fokus;}Sebar dengan dua jari: Tatal halaman penuh atas/bawah/kiri/kanan;

\textbf{Luncur atas/bawah/kiri/kanan dua jari: Tatal;}

\begin{enumerate}
  \item Ketuk tiga kali satu jari: Zum masuk atau pilih (berbeza dengan tiga klik biasa).
  \item Tetapan peribadi
  \item : [VoiceOver] > [Ucapan], boleh melaraskan
  \item kelajuan pembacaan
  \item (lalai kira-kira 50\%, boleh dilaraskan ke 70-80\% untuk pendengaran lebih lancar),
\end{enumerate}

\textbf{nada}, tukar pakej suara berbeza. iOS terbina dalam “Suara Dipertingkat” dan “Pengawal Suara”, yang pertama boleh mengurangkan bunyi pernafasan dan jeda, pendengaran lebih semula jadi.\textbf{Android TalkBack}Laluan pengaktifan\textbf{2. Besarkan Fon Kecil, Jadikan Kabur Jelas — Pelarasan Visual}Gerak isyarat asas

\end{tcbinfo}

\begin{tcbinfo}{}
\textbf{Luncur kiri-kanan satu jari: Tukar elemen fokus;}Ketuk dua kali satu jari: Aktifkan elemen;

\textbf{Luncur atas/bawah/kiri/kanan dua jari: Tatal;}

\begin{enumerate}
  \item Ketuk sekali empat jari: Jeda/sambung pembacaan.
  \item Tetapan peribadi
  \item : [TalkBack] > [Tetapan Teks-ke-Ucapan], boleh melaraskan
  \item kelajuan ucapan
  \item ,
\end{enumerate}

\textbf{nada}, tukar\textbf{enjin ucapan}AI. AI tanpa peranan ialah seorang “ahli serba boleh tetapi tidak mahir dalam apa-apa”; dengan peranan, ia boleh “memasuki” identiti itu, berfikir seperti profesional tersebut.\textbf{2. Besarkan Fon Kecil, Jadikan Kabur Jelas — Pelarasan Visual}Gerak Isyarat Zum\textbf{iOS}: [Tetapan] > [Fungsi Pembantu] > [Zum], selepas diaktifkan,

\end{tcbinfo}

\booksubsection{Besarkan Fon Kecil, Jadikan Kabur Jelas — Pelarasan Visual}

\booksubsection{Gerak Isyarat Zum}

\begin{itemize}
  \item \textbf{Android}, ketuk dua kali tiga jari untuk mengunci nisbah zum semasa.\textbf{Android}: [Tetapan] > [Fungsi Pembantu] > [Zum], ketik tiga jari untuk zum masuk, ketik tiga jari lagi untuk keluar; atau terus\textbf{ketik skrin tiga kali dengan tiga jari}untuk mengaktifkan zum sementara (maksimum 8 kali ganda).
  \item \textbf{iOS}iOS\textbf{: [Tetapan] > [Fungsi Pembantu] > [Paparan \& Saiz Teks], aktifkan suis "Teks Tebal", "Tingkatkan Kontras", "Kurangkan Kelegapan" dan lain-lain. "Songsang Pintar" boleh mengekalkan fon hitam pada latar belakang cerah, mengelakkan fon putih atas latar putih.}Android
\end{itemize}

\booksubsection{Teks Kontras Tinggi / Teks Tebal}

\begin{itemize}
  \item \textbf{Android}iOS
  \item \textbf{iOS}empat mod buta warna merah/hijau/biru/kuning
\end{itemize}

\booksubsection{Penapis Warna (Mesra Buta Warna)}

\begin{itemize}
  \item \textbf{Android}: [Tetapan] > [Fungsi Pembantu] > [Pembetulan Warna], mengandungi\textbf{tiga mod buta warna}+ suis "Songsangkan warna".
  \item \textbf{iOS}Apabila pembaca skrin menemui imej, amalan tradisional adalah melangkau atau hanya membaca perkataan "imej". Tetapi kini AI boleh\textbf{menerangkan kandungan imej dalam satu ayat}, sangat membantu untuk mengenalpasti kod pengesahan, membaca imej, dan mengenalpasti menu.
\end{itemize}

\booksubsection{Penerangan Imej AI: Biar AI Menjadi “Mata” Anda}

iOS 15+\textbf{: Dalam [Tetapan] > [Fungsi Pembantu] > [VoiceOver] > [Penerangan Audio], aktifkan “Penerangan Imej”. Sistem akan memanggil}model AI tempatan

\begin{tcbinfo}{untuk menjana penerangan teks imej secara automatik (memerlukan iPhone 12 dan ke atas).}
\begin{itemize}
  \item \textbf{Aplikasi TalkBack versi baharu}: Google TalkBack dalam versi lebih baharu menyepadukan fungsi\textbf{soal jawab imej}, boleh bertanya "Apa yang ada dalam gambar ini", AI menjawab secara masa nyata (tidak bergantung ketat pada versi sistem Android, kemas kini aplikasi TalkBack sahaja untuk mengalaminya).
  \item \textbf{App pihak ketiga}: "Seeing AI" (iOS), "Lookout" (Android), "Xiao Yi Bang Bang" (Huawei) semuanya adalah pembantu visual AI yang cemerlang, boleh mengenalpasti barang, membaca dokumen, mengenalpasti ekspresi muka.\textbf{Senario Penggunaan}Cuba tanya tentang foto jalanan "Ada berapa orang dalam gambar ini? Apa yang mereka lakukan?", atau buka App penghantaran makanan, biar AI membaca nama hidangan dalam imej menu —
  \item \textbf{ini adalah "teknologi menjadikan kehidupan lebih adil" yang sebenar}.
\end{itemize}

\end{tcbinfo}

\begin{tcbpractice}{4. "Telinga" Pengguna Pekak — Panduan Lengkap Sarikata Masa Nyata}
Sarikata masa nyata (Live Captions) adalah fungsi\textbf{ucapan-ke-teks peringkat sistem}, sangat berguna bagi pengguna kurang pendengaran dalam persekitaran bising, perpustakaan dan senario bahasa asing.

\end{tcbpractice}

\booksubsection{“Telinga” Pengguna Pekak — Panduan Lengkap Sarikata Masa Nyata}

iOS\textbf{: [Tetapan] > [Fungsi Pembantu] > [Sarikata Masa Nyata], buka suis. Sistem akan mengesan semua audio peranti (panggilan, video, podcast, FaceTime, memo suara), memaparkan sarikata secara automatik di bahagian atas skrin.}Android 12+

\begin{tcbinfo}{: Tekan}
\textbf{Android}Klik ikon “Sarikata Masa Nyata” dalam kotak timbul; atau aktifkan dalam [Tetapan] > [Fungsi Pembantu] > [Sarikata Masa Nyata].

\textbf{Tip pengoptimuman ketepatan}:\textbf{Gunakan sebanyak mungkin dalam}persekitaran senyap

\textbf{, jauh dari penyaman udara, dengungan mesin;}

\begin{enumerate}
  \item menyebut dengan jelas dan kelajuan sederhana\textbf{;}Perbualan berbilang orang akan mengurangkan ketepatan, boleh menggunakan mod "hanya paparkan satu orang";
  \item Menyokong bahasa utama seperti Mandarin, Kantonis, Inggeris,\textbf{kadar pengenalan bahasa kecil lebih rendah};
  \item Muzik latar belakang atau gangguan suara manusia semasa bercakap akan menyebabkan pengenalan terputus-putus.
  \item 5. Dwi Saluran Senario Panggilan: Ucapan ke Teks + Teks ke Ucapan\textbf{Apabila pengguna}tidak boleh mendengar dan juga tidak boleh bercakap
  \item (seperti kehilangan suara secara tiba-tiba, tempoh pemulihan pembedahan tekak, gagap teruk), telefon masih boleh menjadi "perantara mesej".
\end{enumerate}

\end{tcbinfo}

\booksubsection{Dwi Saluran Senario Panggilan: Ucapan ke Teks + Teks ke Ucapan}

iOS\textbf{: [Tetapan] > [Fungsi Pembantu] > [RTT/TTY], selepas diaktifkan, papan kekunci akan dipaparkan pada skrin semasa panggilan, suara pihak lain akan ditaip menjadi teks secara masa nyata,}teks yang anda taip juga akan dibacakan

\begin{tcbinfo}{kepada pihak lain.}
\textbf{Android}: [Tetapan] > [Fungsi Pembantu] > [Transliterasi Masa Nyata] + [Pemberitahuan Bunyi]. “Transliterasi Masa Nyata” memaparkan ucapan pihak lain sebagai teks secara masa nyata; “Pemberitahuan Bunyi” menukar input anda kepada ucapan apabila suara tidak mudah.\textbf{Keserasian}:

\textbf{iOS}: Kedua-dua pihak perlu menyokong, teks dipaparkan semasa menaip tanpa kependaman, pengendali utama telah menyokong;

\textbf{Keberserasian}:

\begin{itemize}
  \item \textbf{Fungsi}ucapan-ke-teks
  \item \textbf{WeChat/QQ juga merupakan alternatif: Tekan lama mesej suara "Tukar ke teks".}6. Fungsi "Kecil Tetapi Berguna" untuk Semua Orang
  \item Ketuk belakang\textbf{(iOS 14+): [Tetapan] > [Fungsi Pembantu] > [Sentuh] > [Ketuk Belakang], boleh mengikatkan}ketukan dua kali/tiga kali
\end{itemize}

\end{tcbinfo}

\booksubsection{Fungsi “Kecil Tetapi Berguna” untuk Semua Orang}

\begin{itemize}
  \item \textbf{Sentuh pembantu / Bola terapung}: iOS [Fungsi Pembantu] > [Sentuh] > [Sentuh Pembantu], Android [Fungsi Pembantu] > [Navigasi Terapung]. Butang maya menggantikan gerak isyarat atau kekunci fizikal, menyelamatkan keadaan apabila kekunci tidak berfungsi atau pengendalian satu tangan.\textbf{Kawalan suis}: Mengawal telefon melalui suis luaran, klik tunggal, pergerakan kepala,
  \item \textbf{direka untuk pengguna kecacatan anggota badan}. iOS [Fungsi Pembantu] > [Kawalan Suis], Android [Fungsi Pembantu] > [Kawalan Suis].
  \item \textbf{Pengenalan bunyi}(iOS): [Fungsi Pembantu] > [Pengenalan Bunyi], selepas diaktifkan, telefon akan "mendengar" loceng pintu, tangisan bayi, penggera kebakaran, siren dan lain-lain bunyi,\textbf{menggetar + pemberitahuan skrin apabila dikenalpasti}, sangat berharga bagi pengguna pekak dan semasa tidur.
  \item \textbf{Peringatan kilat LED}: iOS [Fungsi Pembantu] > [Audio/Visual] > [Kilat LED untuk amaran], Android [Tetapan] > [Pemberitahuan] > [Peringatan lampu kilat]. Lampu kilat berkelip semasa panggilan masuk dan pemberitahuan,\textbf{boleh dirasai tanpa bunyi dalam persekitaran bising}.
  \item \textbf{Audio mono}: [Tetapan] > [Fungsi Pembantu] > [Audio] > [Audio Mono], menggabungkan saluran kiri dan kanan. Pengguna yang mempunyai gangguan pendengaran satu telinga tidak akan kehilangan separuh maklumat bunyi.
  \item \textbf{Falsafah Reka Bentuk Aksesibiliti Apple}: Apple melabur sumber yang sangat besar dalam bidang aksesibiliti, mengemas kini dan menambah fungsi aksesibiliti setiap tahun. Laman web rasmi Apple mempunyai halaman aksesibiliti khusus yang memperincikan kaedah penggunaan setiap fungsi. Fungsi aksesibiliti iOS biasanya lebih lengkap daripada Android dengan pengalaman yang lebih lancar, dan itu sebab penting ramai pengguna kurang upaya memilih Apple.
  \item \textbf{Luangkan 30 Saat untuk Mencubanya}Cari masa lapang, lakukan pengalaman berikut:
\end{itemize}

\begin{tcbinfo}{Aktifkan}
VoiceOver / TalkBack

\end{tcbinfo}

\begin{tcbpractice}{, cuba melayari WeChat dengan luncuran tiga jari tanpa menggunakan mata;}
Buka

\begin{enumerate}
  \item Sarikata Masa Nyata\textbf{, bercakap ke telefon dan lihat ketepatan pengenalan;}Gunakan
  \item ketikan tiga jari\textbf{untuk membesarkan teks hingga 200\%, alami kehidupan harian pengguna celik;}Matikan semua fungsi sebelum ini, fikirkan "Jika saya sentiasa memerlukan fungsi ini, bagaimana saya mahu ia direka bentuk".
  \item Bab 6: Lindungi Kehidupan Digital Anda — Keselamatan \& Privasi\textbf{Keselamatan telefon = keselamatan privasi + keselamatan harta. Dalam era data sebagai aset ini, melindungi telefon bermaksud melindungi diri anda.}Satu, Kunci Skrin: Garis Pertahanan Fizikal Pertama
  \item Kunci skrin adalah garis pertahanan pertama telefon anda, dan juga yang paling mudah diabaikan.
\end{enumerate}

\end{tcbpractice}
